\begin{figure*}[t]
    \centering
    \includegraphics[width=\linewidth]{figures/fig3_zeroshot.pdf}
    \caption{\textbf{Zero-shot cross-embodiment generalization performance.} LAP-3B achieves performance comparable to the $\pi_{0.5}$-DROID on the seen embodiment. Across three previously unseen embodiments and six real-world manipulation tasks, LAP-3B attains over \textbf{50\% average zero-shot success}, delivering approximately a $2\times$ improvement over the strongest baselines, while all open-sourced VLAs collapse to zero success rate. Error bars denote 95\% finite-sample–valid confidence intervals that control the Type-I error (miscoverage) probability, following~\cite{vincent2024generalizablebehaviorcloningpolicy}.
}
    \label{fig:zero-shot}
\end{figure*}

\section{Experiments}
\label{sec:experiments}

Our experiments aim to evaluate whether language-actions provide a superior action representation for VLA pre-training. In particular, we focus on \textbf{cross-embodiment generalization}, which remains a significant limitation of state-of-the-art VLAs. Concretely, we seek to answer:
{
% \renewcommand{\theenumi}{\Alph{enumi}}
% \renewcommand{\labelenumi}{\theenumi.}
\begin{enumerate}
    \item Does LAP enable zero-shot cross-embodiment transfer?
    \item Does LAP support efficient fine-tuning to new embodiments on challenging tasks?
    \item How does LAP’s language-action representation shape the learned representations and enable strong cross-embodiment generalization?
    \item What is the effect of co-training \textsc{LAP-3B} with VQA datasets?
    \item Does LAP scale favorably with model size?
\end{enumerate}
}
To address these questions, we evaluate LAP across \textbf{four robot embodiments}, \textbf{ten real-world manipulation tasks}, and the LIBERO simulation benchmark \cite{liu2023liberobenchmarkingknowledgetransfer}. We first test out-of-the-box performance on one seen and three unseen embodiments, then study fine-tuning efficiency on new robots and tasks. We investigate the impact of LAP's language-action representation on training dynamics and conclude with an analysis of co-training with Visual--Question--Answering (VQA) datasets.

\noindent\textbf{Baselines.}
We evaluate two categories of baselines below. All replicated baselines are carefully tuned (learning rate, batch size, loss weights, etc.) for fair comparison; full details are in Appendix~\ref{app:implementation_details}.


\paragraph{Open-Sourced VLAs.}
We evaluate several publicly released VLA checkpoints to assess whether existing methods already exhibit zero-shot cross-embodiment transfer:
(1) \textbf{\bm{$\pi_{0.5}$}-DROID} \cite{intelligence2025pi05visionlanguageactionmodelopenworld}, currently the strongest VLA on the DROID benchmark;
(2) \textbf{\bm{$\pi_{0.5}$}-Base} \cite{intelligence2025pi05visionlanguageactionmodelopenworld}, the strongest publicly available VLA base model to our knowledge;
(3) \textbf{X-VLA} \cite{zheng2025xvlasoftpromptedtransformerscalable}, which explicitly targets embodiment generalization;
(4) \textbf{MolmoAct} \cite{lee2025molmoactactionreasoningmodels}; and
(5) \textbf{OpenVLA} \cite{kim2024openvlaopensourcevisionlanguageactionmodel}.


\paragraph{Action-Representation Comparison.}
To isolate the impact of language-action supervision, we train controlled baselines with identical architectures and data mixtures, differing only in action representation:

\begin{enumerate}
    \item \textbf{$\bm{\pi_{0.5}}$-replicated}: supervises the VLM using FAST tokens \cite{pertsch2025fastefficientactiontokenization}. Since the FAST tokenizer is pre-trained on substantially more robot data than our model, this forms a particularly strong baseline.
    \item \textbf{$\bm{\pi_{0}}$-replicated}: removes action supervision from the VLM entirely and only supervises the action expert, measuring the benefit of language-action pre-training.
    \item \textbf{VLA-0-replicated}: supervises the VLM following \citet{goyal2025vla0buildingstateoftheartvlas}, where continuous actions are represented as digit tokens.
\end{enumerate}
Across all experiments, we conduct \textbf{over 1300 real-robot evaluation trials}.

\subsection{Zero-Shot Cross-Embodiment Generalization}
\label{subsec:zero_shot}



\noindent\textbf{Embodiments and Tasks.}
We evaluate on four robot embodiments (Fig.~\ref{fig:zero-shot}): the DROID setup \cite{khazatsky2025droidlargescaleinthewildrobot} which is included in the pre-training mixture, three unseen embodiments---a custom 7-DoF Franka Panda robot with a novel gripper and wrist-mounted camera configuration, the 6-DoF YAM robot, and the 7-DoF Kinova robot.

We evaluate five manipulation task types (Fig.~\ref{fig:zero-shot}, top row):
\begin{enumerate}
    \item \textbf{Pick and Place}: pick up an object and place it into a container, testing generalization across gripper geometry, camera placement, and robot appearance.
    \item \textbf{Sort}: sort all objects on the table into a basket, evaluating long-horizon manipulation.
    \item \textbf{Tissue}: pull tissue from a box and place it on the table, requiring adaptive grasping of deformable objects and height reasoning.
    \item \textbf{Put Towel}: pick up a flat towel and place it into a basket, requiring precise grasp localization and large vertical motion.
    \item \textbf{Pour}: pour contents from a bowl, testing rotational control and primitive sequencing.
\end{enumerate}

These tasks require full 6-DoF manipulation, e.g., rotating the gripper to avoid collisions during sorting, and probe complementary primitive skills.

For each embodiment, we evaluate on two of these tasks with 20 trials per task. We report mean success rates with 95\% confidence intervals following the method of \cite{vincent2024generalizablebehaviorcloningpolicy}, which provides statistically valid finite-sample coverage. Unlike heuristic standard deviations, these intervals explicitly control the Type-I error (miscoverage) probability. Success metric is binary: a trial is marked successful only if the full task is completed. Results are summarized in Fig.~\ref{fig:zero-shot}.

\noindent\textbf{Seen Embodiment Evaluation.} On the in-distribution DROID setup, we observe that nearly all open-sourced VLAs fail to achieve non-trivial task success, with the sole exception of $\pi_{0.5}-\text{DROID}$ \cite{intelligence2025pi05visionlanguageactionmodelopenworld}, which has been explicitly fine-tuned on the DROID dataset \cite{khazatsky2025droidlargescaleinthewildrobot}. This highlights a critical limitation of current VLAs: despite large-scale multi-embodiment pre-training, effective deployment even on seen embodiments typically requires embodiment-specific adaptation. 

For our replicated comparison baselines, $\pi_{0.5}$-replicated and $\pi_{0}$-replicated achieve moderate and comparable performance. However, $\text{VLA-0}$-replicated consistently collapses to near-zero actions. We find this failure mode stems from overfitting to trivial predictions, which manifests as predicting only small-magnitude action tokens. This issue may be less of a concern on smaller, cleaner benchmarks such as LIBERO \cite{liu2023liberobenchmarkingknowledgetransfer} but becomes detrimental at large scale, where noise and behavioral diversity are unavoidable.

Notably, \textsc{LAP-3B} consistently outperforms $\pi_{0.5}$-replicated and $\pi_{0}$-replicated by approximately 15 percentage points across tasks, despite sharing the same backbone architecture and training data; this highlights the superiority of the language-action representation. Moreover, \textsc{LAP-3B} achieves performance on par with $\pi_{0.5}$-DROID---without any embodiment-specific fine-tuning on the DROID dataset itself. This indicates that language-action supervision enables effective learning from large-scale, heterogeneous robot data, leading to strong in-distribution performance.


\noindent\textbf{Unseen Embodiment Evaluation.} We next evaluate zero-shot transfer to three previously unseen robot embodiments. Across all settings, \textsc{LAP-3B} demonstrates a substantial and consistent advantage, with an \emph{average success rate exceeding 50\%}, representing an approximate \textbf{2$\times$ improvement over the strongest baseline}. Importantly, \textsc{LAP-3B} is the only model to achieve consistently high performance across all novel embodiments.

In contrast, all open-sourced VLAs, including $\pi_{0.5}$-DROID, which performed strongly in-distribution, fail to generate meaningful behavior on unseen robots. This sharp collapse suggests that existing VLA training recipes remain tightly coupled to embodiment-specific control conventions, preventing transfer even when visual and task semantics remain similar.


By contrast, \textsc{LAP-3B} exhibits both \emph{semantically grounded movements} and \emph{fine-grained spatial control} across embodiments. We hypothesize that the language-action interface provides a shared, embodiment-agnostic abstraction that better preserves the VLM’s generalizable visual-semantic representations while enabling precise execution under novel robot kinematics. This decoupling allows \textsc{LAP-3B} to retain coherent task intent while avoiding the embodiment-specific overfitting that hampers existing VLAs.

\begin{figure*}[t]
    \centering
    \includegraphics[width=\linewidth]{figures/fig4_finetune.pdf}
    \caption{\textbf{Fine-tuning efficiency in simulation (LIBERO) and real-world manipulation tasks}. Across both domains, \textsc{LAP-3B} adapts substantially faster than baseline policies, reaching high performance with significantly fewer epochs and demonstrations. In simulation, \textsc{LAP-3B} converges to near-optimal success within a fraction of the training steps required by baselines. On real robots, \textsc{LAP-3B} achieves comparable task performance using approximately \textbf{$\bm{2.5\times}$ fewer demonstrations}, demonstrating substantially improved data and compute efficiency when transferring to new embodiments.}
    \label{fig:libero-fine-tuning}
\end{figure*}
\vspace{2pt}

